% ============================================================
% Section 69: 二维密排结构的布里渊区与紧束缚能带
% ============================================================
\section{固体理论：二维密排结构的布里渊区与紧束缚能带}
\label{sec:2d-triangular-lattice}

\begin{modelbox}[题目：二维密排结构的布里渊区、费米波矢与$s$带]
考虑原子间距为 $a$ 的由同种原子构成的二维密排结构。
\begin{enumerate}
    \item 画出第一和第二布里渊区；
    \item 求出每个原子有一个自由电子时的费米波矢；
    \item 利用紧束缚近似，推导$s$态原子波函数形成的能带。
\end{enumerate}
\end{modelbox}

\begin{tipbox}[考点快速分析]
\begin{itemize}
    \item \textbf{二维密排}：三角晶格，配位数 6，原胞面积 $\frac{\sqrt{3}}{2}a^2$
    \item \textbf{倒格子}：也是三角晶格，倒格常数 $\frac{4\pi}{\sqrt{3}a}$
    \item \textbf{布里渊区}：第一区为正六边形；第二区为外围 6 个三角形
    \item \textbf{自由电子}：二维态密度与费米圆面积相关
    \item \textbf{紧束缚}：6 个最近邻贡献，三角恒等式化简
\end{itemize}
\end{tipbox}

\begin{derivationbox}[标准解题过程]
\textbf{Step 1: 第一和第二布里渊区}

二维密排结构即\textbf{三角晶格}。取正格子原胞基矢（见图示）：
\begin{equation}
\bm{a}_1 = a\,\hat{\bm{i}},\qquad
\bm{a}_2 = \frac{a}{2}\,\hat{\bm{i}} + \frac{\sqrt{3}a}{2}\,\hat{\bm{j}}.
\end{equation}

由 $\bm{a}_i\cdot\bm{b}_j = 2\pi\delta_{ij}$，求得倒格子基矢：
\begin{equation}
\bm{b}_1 = \frac{4\pi}{\sqrt{3}a}\biggl(\frac{\sqrt{3}}{2}\,\hat{\bm{i}} - \frac{1}{2}\,\hat{\bm{j}}\biggr),\qquad
\bm{b}_2 = \frac{4\pi}{\sqrt{3}a}\,\hat{\bm{j}}.
\end{equation}

倒格子也是三角晶格。选取一倒格点为原点，作到最近邻 6 个倒格点（$\pm\bm{b}_1$、$\pm\bm{b}_2$、$\pm(\bm{b}_1+\bm{b}_2)$）连线的中垂线。

\begin{itemize}
    \item \textbf{第一布里渊区}：围绕原点的\textbf{正六边形}。
    \item \textbf{第二布里渊区}：第一区外由中垂线围成的\textbf{6 个三角形}区域。
\end{itemize}

\textbf{Step 2: 每个原子一个自由电子时的费米波矢}

设晶体面积为 $S$，原子总数为 $N$。每个原子贡献 1 个自由电子，总电子数 $N_e=N$。

二维自由电子模型中，考虑自旋简并（因子 2），$k$ 空间态密度为 $S/(2\pi)^2$。费米圆面积 $\pi k_F^2$ 内的状态数等于总电子数：
\begin{equation}
2\cdot\frac{S}{(2\pi)^2}\cdot\pi k_F^2 = N.
\end{equation}
化简得
\begin{equation}
k_F^2 = 2\pi\,\frac{N}{S} = 2\pi n,
\qquad\text{即}\qquad
k_F = \sqrt{2\pi n},
\end{equation}
其中 $n=N/S$ 为电子面密度。

三角晶格一个原胞面积为 $|\bm{a}_1\times\bm{a}_2| = \frac{\sqrt{3}}{2}a^2$，含 1 个原子。故
\begin{equation}
n = \frac{2}{\sqrt{3}a^2}.
\end{equation}
代入得
\begin{equation}
\boxed{k_F = \sqrt{\frac{4\pi}{\sqrt{3}a^2}} = \frac{2}{a}\sqrt{\frac{\pi}{\sqrt{3}}} \approx \frac{1.70}{a}.}
\end{equation}
数值上 $k_F \approx 2/a\times\sqrt{\pi/1.732} \approx 2/a\times\sqrt{1.814} \approx 2/a\times1.347 \approx 2.69/a$。

精确值：$k_F = \frac{2}{a}\bigl(\frac{\pi}{\sqrt{3}}\bigr)^{1/2} \approx 2.69/a$。

\textbf{Step 3: 紧束缚近似下的$s$带}

以任一原子为原点，6 个最近邻的相对坐标为
\begin{equation}
(a,0),\; (-a,0),\; \Bigl(\frac{a}{2},\frac{\sqrt{3}a}{2}\Bigr),\; \Bigl(-\frac{a}{2},-\frac{\sqrt{3}a}{2}\Bigr),\; \Bigl(\frac{a}{2},-\frac{\sqrt{3}a}{2}\Bigr),\; \Bigl(-\frac{a}{2},\frac{\sqrt{3}a}{2}\Bigr).
\end{equation}

紧束缚能带公式（$s$态球对称，6 个方向 $J_s$ 相同）：
\begin{equation}
E(\bm{k}) = E_s + J_{0,0}^{s} + J_s\sum_{\bm{R}_l}e^{i\bm{k}\cdot\bm{R}_l}.
\end{equation}

代入 6 个最近邻：
\begin{align}
E(\bm{k}) &= E_s + J_{0,0}^{s} + J_s\Bigl[e^{ik_xa} + e^{-ik_xa}
+ e^{i(\frac{a}{2}k_x + \frac{\sqrt{3}a}{2}k_y)} + e^{-i(\frac{a}{2}k_x + \frac{\sqrt{3}a}{2}k_y)} \notag\\
&\qquad + e^{i(\frac{a}{2}k_x - \frac{\sqrt{3}a}{2}k_y)} + e^{-i(\frac{a}{2}k_x - \frac{\sqrt{3}a}{2}k_y)}\Bigr] \notag\\
&= E_s + J_{0,0}^{s} + 2J_s\Bigl[\cos(k_xa)
+ \cos\Bigl(\frac{ak_x}{2} + \frac{\sqrt{3}ak_y}{2}\Bigr)
+ \cos\Bigl(\frac{ak_x}{2} - \frac{\sqrt{3}ak_y}{2}\Bigr)\Bigr].
\end{align}

利用 $\cos(\alpha+\beta)+\cos(\alpha-\beta)=2\cos\alpha\cos\beta$，令 $\alpha=ak_x/2$，$\beta=\sqrt{3}ak_y/2$：
\begin{equation}
\boxed{E(\bm{k}) = E_s + J_{0,0}^{s} + 2J_s\Bigl[\cos(k_xa) + 2\cos\Bigl(\frac{ak_x}{2}\Bigr)\cos\Bigl(\frac{\sqrt{3}ak_y}{2}\Bigr)\Bigr].}
\end{equation}
\end{derivationbox}

\begin{formulabox}[答案汇总]
\begin{center}
\begin{tabular}{ll}
\toprule
\textbf{问题} & \textbf{答案} \\
\midrule
(1) 第一布里渊区 & 正六边形 \\
(1) 第二布里渊区 & 外围 6 个三角形 \\
(2) 费米波矢 & $k_F = \sqrt{4\pi/(\sqrt{3}a^2)} \approx 2.69/a$ \\
(3) $s$带能带 & $E(\bm{k}) = E_s+J_{0,0}^{s}+2J_s\bigl[\cos(k_xa)+2\cos\frac{ak_x}{2}\cos\frac{\sqrt{3}ak_y}{2}\bigr]$ \\
\bottomrule
\end{tabular}
\end{center}
\end{formulabox}

\begin{insightbox}[物理图像]
\begin{itemize}
    \item 二维密排结构的布里渊区为正六边形，与石墨烯的布里渊区形状相同（但石墨烯是复式格子，有额外的能带结构）。
    \item 自由电子费米波矢仅由电子面密度决定，与晶格类型无关（仅通过原子密度间接关联）。
    \item 三角晶格有 6 个最近邻，能带公式中出现三项余弦叠加，在特殊高对称方向可进一步化简。
\end{itemize}
\end{insightbox}
